\subsection*{Вариант 1 для  Каниной Агнии. Дроби \\ \today   \currenttime} 
  \begin{enumerate}
\item Найди число, если его  19 \%  равна 11400. Нарисуй схему.  \\ \nopagebreak \begin{tikzpicture}
 \draw[step=0.7cm ,gray,very thin] (0,0) grid (18,6);
 % Обводка внешнего контура  \draw[black, thick] (0,0) rectangle (18,6);
 \end{tikzpicture}
\item Найди число, если его $\frac{1}{6}$  часть равна 60. Нарисуй схему.  \\ \nopagebreak \begin{tikzpicture}
 \draw[step=0.7cm ,gray,very thin] (0,0) grid (18,6);
 % Обводка внешнего контура  \draw[black, thick] (0,0) rectangle (18,6);
 \end{tikzpicture}
\item Найди число, если его  18 \%  равна 28800. Нарисуй схему.  \\ \nopagebreak \begin{tikzpicture}
 \draw[step=0.7cm ,gray,very thin] (0,0) grid (18,6);
 % Обводка внешнего контура  \draw[black, thick] (0,0) rectangle (18,6);
 \end{tikzpicture}
\item Найди число, если его $\frac{4}{15}$  часть равна 60. Нарисуй схему.  \\ \nopagebreak \begin{tikzpicture}
 \draw[step=0.7cm ,gray,very thin] (0,0) grid (18,6);
 % Обводка внешнего контура  \draw[black, thick] (0,0) rectangle (18,6);
 \end{tikzpicture}
\item Найди число, если его $\frac{9}{14}$  часть равна 126. Нарисуй схему.  \\ \nopagebreak \begin{tikzpicture}
 \draw[step=0.7cm ,gray,very thin] (0,0) grid (18,6);
 % Обводка внешнего контура  \draw[black, thick] (0,0) rectangle (18,6);
 \end{tikzpicture}
\item Найди число, если его $\frac{1}{8}$  часть равна 32. Нарисуй схему.  \\ \nopagebreak \begin{tikzpicture}
 \draw[step=0.7cm ,gray,very thin] (0,0) grid (18,6);
 % Обводка внешнего контура  \draw[black, thick] (0,0) rectangle (18,6);
 \end{tikzpicture}
\end{enumerate}  
\newpage
\subsection*{Вариант 2 для  Спицына Марка. Дроби \\ \today   \currenttime} 
  \begin{enumerate}
\item Найди число, если его $\frac{1}{3}$  часть равна 15. Нарисуй схему.  \\ \nopagebreak \begin{tikzpicture}
 \draw[step=0.7cm ,gray,very thin] (0,0) grid (18,6);
 % Обводка внешнего контура  \draw[black, thick] (0,0) rectangle (18,6);
 \end{tikzpicture}
\item Найди число, если его  7 \%  равна 7700. Нарисуй схему.  \\ \nopagebreak \begin{tikzpicture}
 \draw[step=0.7cm ,gray,very thin] (0,0) grid (18,6);
 % Обводка внешнего контура  \draw[black, thick] (0,0) rectangle (18,6);
 \end{tikzpicture}
\item Найди число, если его $\frac{3}{10}$  часть равна 30. Нарисуй схему.  \\ \nopagebreak \begin{tikzpicture}
 \draw[step=0.7cm ,gray,very thin] (0,0) grid (18,6);
 % Обводка внешнего контура  \draw[black, thick] (0,0) rectangle (18,6);
 \end{tikzpicture}
\item Найди число, если его  15 \%  равна 6000. Нарисуй схему.  \\ \nopagebreak \begin{tikzpicture}
 \draw[step=0.7cm ,gray,very thin] (0,0) grid (18,6);
 % Обводка внешнего контура  \draw[black, thick] (0,0) rectangle (18,6);
 \end{tikzpicture}
\item Найди число, если его $\frac{1}{4}$  часть равна 24. Нарисуй схему.  \\ \nopagebreak \begin{tikzpicture}
 \draw[step=0.7cm ,gray,very thin] (0,0) grid (18,6);
 % Обводка внешнего контура  \draw[black, thick] (0,0) rectangle (18,6);
 \end{tikzpicture}
\item Найди число, если его $\frac{2}{10}$  часть равна 20. Нарисуй схему.  \\ \nopagebreak \begin{tikzpicture}
 \draw[step=0.7cm ,gray,very thin] (0,0) grid (18,6);
 % Обводка внешнего контура  \draw[black, thick] (0,0) rectangle (18,6);
 \end{tikzpicture}
\end{enumerate}  
\newpage
\subsection*{Вариант 3 для  Александровой Ксении Андреевны. Дроби \\ \today   \currenttime} 
  \begin{enumerate}
\item Найди число, если его $\frac{1}{4}$  часть равна 32. Нарисуй схему.  \\ \nopagebreak \begin{tikzpicture}
 \draw[step=0.7cm ,gray,very thin] (0,0) grid (18,6);
 % Обводка внешнего контура  \draw[black, thick] (0,0) rectangle (18,6);
 \end{tikzpicture}
\item Найди число, если его $\frac{1}{4}$  часть равна 36. Нарисуй схему.  \\ \nopagebreak \begin{tikzpicture}
 \draw[step=0.7cm ,gray,very thin] (0,0) grid (18,6);
 % Обводка внешнего контура  \draw[black, thick] (0,0) rectangle (18,6);
 \end{tikzpicture}
\item Найди число, если его $\frac{5}{9}$  часть равна 45. Нарисуй схему.  \\ \nopagebreak \begin{tikzpicture}
 \draw[step=0.7cm ,gray,very thin] (0,0) grid (18,6);
 % Обводка внешнего контура  \draw[black, thick] (0,0) rectangle (18,6);
 \end{tikzpicture}
\item Найди число, если его $\frac{7}{17}$  часть равна 119. Нарисуй схему.  \\ \nopagebreak \begin{tikzpicture}
 \draw[step=0.7cm ,gray,very thin] (0,0) grid (18,6);
 % Обводка внешнего контура  \draw[black, thick] (0,0) rectangle (18,6);
 \end{tikzpicture}
\item Найди число, если его  16 \%  равна 6400. Нарисуй схему.  \\ \nopagebreak \begin{tikzpicture}
 \draw[step=0.7cm ,gray,very thin] (0,0) grid (18,6);
 % Обводка внешнего контура  \draw[black, thick] (0,0) rectangle (18,6);
 \end{tikzpicture}
\item Найди число, если его  7 \%  равна 6300. Нарисуй схему.  \\ \nopagebreak \begin{tikzpicture}
 \draw[step=0.7cm ,gray,very thin] (0,0) grid (18,6);
 % Обводка внешнего контура  \draw[black, thick] (0,0) rectangle (18,6);
 \end{tikzpicture}
\end{enumerate}  
\newpage
\subsection*{Вариант 4 для  Баумгертнер Елизаветы Владиславовны. Дроби \\ \today   \currenttime} 
  \begin{enumerate}
\item Найди число, если его  13 \%  равна 18200. Нарисуй схему.  \\ \nopagebreak \begin{tikzpicture}
 \draw[step=0.7cm ,gray,very thin] (0,0) grid (18,6);
 % Обводка внешнего контура  \draw[black, thick] (0,0) rectangle (18,6);
 \end{tikzpicture}
\item Найди число, если его  12 \%  равна 16800. Нарисуй схему.  \\ \nopagebreak \begin{tikzpicture}
 \draw[step=0.7cm ,gray,very thin] (0,0) grid (18,6);
 % Обводка внешнего контура  \draw[black, thick] (0,0) rectangle (18,6);
 \end{tikzpicture}
\item Найди число, если его $\frac{4}{11}$  часть равна 44. Нарисуй схему.  \\ \nopagebreak \begin{tikzpicture}
 \draw[step=0.7cm ,gray,very thin] (0,0) grid (18,6);
 % Обводка внешнего контура  \draw[black, thick] (0,0) rectangle (18,6);
 \end{tikzpicture}
\item Найди число, если его $\frac{1}{4}$  часть равна 8. Нарисуй схему.  \\ \nopagebreak \begin{tikzpicture}
 \draw[step=0.7cm ,gray,very thin] (0,0) grid (18,6);
 % Обводка внешнего контура  \draw[black, thick] (0,0) rectangle (18,6);
 \end{tikzpicture}
\item Найди число, если его $\frac{7}{12}$  часть равна 84. Нарисуй схему.  \\ \nopagebreak \begin{tikzpicture}
 \draw[step=0.7cm ,gray,very thin] (0,0) grid (18,6);
 % Обводка внешнего контура  \draw[black, thick] (0,0) rectangle (18,6);
 \end{tikzpicture}
\item Найди число, если его $\frac{1}{4}$  часть равна 16. Нарисуй схему.  \\ \nopagebreak \begin{tikzpicture}
 \draw[step=0.7cm ,gray,very thin] (0,0) grid (18,6);
 % Обводка внешнего контура  \draw[black, thick] (0,0) rectangle (18,6);
 \end{tikzpicture}
\end{enumerate}  
\newpage
\subsection*{Вариант 5 для  Гайдадиной Анны Алексеевны. Дроби \\ \today   \currenttime} 
  \begin{enumerate}
\item Найди число, если его $\frac{1}{4}$  часть равна 12. Нарисуй схему.  \\ \nopagebreak \begin{tikzpicture}
 \draw[step=0.7cm ,gray,very thin] (0,0) grid (18,6);
 % Обводка внешнего контура  \draw[black, thick] (0,0) rectangle (18,6);
 \end{tikzpicture}
\item Найди число, если его $\frac{2}{7}$  часть равна 14. Нарисуй схему.  \\ \nopagebreak \begin{tikzpicture}
 \draw[step=0.7cm ,gray,very thin] (0,0) grid (18,6);
 % Обводка внешнего контура  \draw[black, thick] (0,0) rectangle (18,6);
 \end{tikzpicture}
\item Найди число, если его $\frac{10}{18}$  часть равна 180. Нарисуй схему.  \\ \nopagebreak \begin{tikzpicture}
 \draw[step=0.7cm ,gray,very thin] (0,0) grid (18,6);
 % Обводка внешнего контура  \draw[black, thick] (0,0) rectangle (18,6);
 \end{tikzpicture}
\item Найди число, если его  19 \%  равна 19000. Нарисуй схему.  \\ \nopagebreak \begin{tikzpicture}
 \draw[step=0.7cm ,gray,very thin] (0,0) grid (18,6);
 % Обводка внешнего контура  \draw[black, thick] (0,0) rectangle (18,6);
 \end{tikzpicture}
\item Найди число, если его  4 \%  равна 2800. Нарисуй схему.  \\ \nopagebreak \begin{tikzpicture}
 \draw[step=0.7cm ,gray,very thin] (0,0) grid (18,6);
 % Обводка внешнего контура  \draw[black, thick] (0,0) rectangle (18,6);
 \end{tikzpicture}
\item Найди число, если его $\frac{1}{10}$  часть равна 100. Нарисуй схему.  \\ \nopagebreak \begin{tikzpicture}
 \draw[step=0.7cm ,gray,very thin] (0,0) grid (18,6);
 % Обводка внешнего контура  \draw[black, thick] (0,0) rectangle (18,6);
 \end{tikzpicture}
\end{enumerate}  
\newpage
\subsection*{Вариант 6 для  Гринчук Арины. Дроби \\ \today   \currenttime} 
  \begin{enumerate}
\item Найди число, если его $\frac{1}{5}$  часть равна 25. Нарисуй схему.  \\ \nopagebreak \begin{tikzpicture}
 \draw[step=0.7cm ,gray,very thin] (0,0) grid (18,6);
 % Обводка внешнего контура  \draw[black, thick] (0,0) rectangle (18,6);
 \end{tikzpicture}
\item Найди число, если его  19 \%  равна 19000. Нарисуй схему.  \\ \nopagebreak \begin{tikzpicture}
 \draw[step=0.7cm ,gray,very thin] (0,0) grid (18,6);
 % Обводка внешнего контура  \draw[black, thick] (0,0) rectangle (18,6);
 \end{tikzpicture}
\item Найди число, если его $\frac{1}{7}$  часть равна 35. Нарисуй схему.  \\ \nopagebreak \begin{tikzpicture}
 \draw[step=0.7cm ,gray,very thin] (0,0) grid (18,6);
 % Обводка внешнего контура  \draw[black, thick] (0,0) rectangle (18,6);
 \end{tikzpicture}
\item Найди число, если его $\frac{5}{11}$  часть равна 55. Нарисуй схему.  \\ \nopagebreak \begin{tikzpicture}
 \draw[step=0.7cm ,gray,very thin] (0,0) grid (18,6);
 % Обводка внешнего контура  \draw[black, thick] (0,0) rectangle (18,6);
 \end{tikzpicture}
\item Найди число, если его $\frac{9}{20}$  часть равна 180. Нарисуй схему.  \\ \nopagebreak \begin{tikzpicture}
 \draw[step=0.7cm ,gray,very thin] (0,0) grid (18,6);
 % Обводка внешнего контура  \draw[black, thick] (0,0) rectangle (18,6);
 \end{tikzpicture}
\item Найди число, если его  9 \%  равна 3600. Нарисуй схему.  \\ \nopagebreak \begin{tikzpicture}
 \draw[step=0.7cm ,gray,very thin] (0,0) grid (18,6);
 % Обводка внешнего контура  \draw[black, thick] (0,0) rectangle (18,6);
 \end{tikzpicture}
\end{enumerate}  
\newpage
\subsection*{Вариант 7 для  Ермолаевой Софии. Дроби \\ \today   \currenttime} 
  \begin{enumerate}
\item Найди число, если его $\frac{6}{9}$  часть равна 54. Нарисуй схему.  \\ \nopagebreak \begin{tikzpicture}
 \draw[step=0.7cm ,gray,very thin] (0,0) grid (18,6);
 % Обводка внешнего контура  \draw[black, thick] (0,0) rectangle (18,6);
 \end{tikzpicture}
\item Найди число, если его $\frac{8}{17}$  часть равна 136. Нарисуй схему.  \\ \nopagebreak \begin{tikzpicture}
 \draw[step=0.7cm ,gray,very thin] (0,0) grid (18,6);
 % Обводка внешнего контура  \draw[black, thick] (0,0) rectangle (18,6);
 \end{tikzpicture}
\item Найди число, если его  13 \%  равна 24700. Нарисуй схему.  \\ \nopagebreak \begin{tikzpicture}
 \draw[step=0.7cm ,gray,very thin] (0,0) grid (18,6);
 % Обводка внешнего контура  \draw[black, thick] (0,0) rectangle (18,6);
 \end{tikzpicture}
\item Найди число, если его  8 \%  равна 10400. Нарисуй схему.  \\ \nopagebreak \begin{tikzpicture}
 \draw[step=0.7cm ,gray,very thin] (0,0) grid (18,6);
 % Обводка внешнего контура  \draw[black, thick] (0,0) rectangle (18,6);
 \end{tikzpicture}
\item Найди число, если его $\frac{1}{7}$  часть равна 21. Нарисуй схему.  \\ \nopagebreak \begin{tikzpicture}
 \draw[step=0.7cm ,gray,very thin] (0,0) grid (18,6);
 % Обводка внешнего контура  \draw[black, thick] (0,0) rectangle (18,6);
 \end{tikzpicture}
\item Найди число, если его $\frac{1}{7}$  часть равна 63. Нарисуй схему.  \\ \nopagebreak \begin{tikzpicture}
 \draw[step=0.7cm ,gray,very thin] (0,0) grid (18,6);
 % Обводка внешнего контура  \draw[black, thick] (0,0) rectangle (18,6);
 \end{tikzpicture}
\end{enumerate}  
\newpage
\subsection*{Вариант 8 для  Икаева Руслана Алановича. Дроби \\ \today   \currenttime} 
  \begin{enumerate}
\item Найди число, если его $\frac{1}{4}$  часть равна 36. Нарисуй схему.  \\ \nopagebreak \begin{tikzpicture}
 \draw[step=0.7cm ,gray,very thin] (0,0) grid (18,6);
 % Обводка внешнего контура  \draw[black, thick] (0,0) rectangle (18,6);
 \end{tikzpicture}
\item Найди число, если его  20 \%  равна 24000. Нарисуй схему.  \\ \nopagebreak \begin{tikzpicture}
 \draw[step=0.7cm ,gray,very thin] (0,0) grid (18,6);
 % Обводка внешнего контура  \draw[black, thick] (0,0) rectangle (18,6);
 \end{tikzpicture}
\item Найди число, если его $\frac{7}{18}$  часть равна 126. Нарисуй схему.  \\ \nopagebreak \begin{tikzpicture}
 \draw[step=0.7cm ,gray,very thin] (0,0) grid (18,6);
 % Обводка внешнего контура  \draw[black, thick] (0,0) rectangle (18,6);
 \end{tikzpicture}
\item Найди число, если его $\frac{7}{19}$  часть равна 133. Нарисуй схему.  \\ \nopagebreak \begin{tikzpicture}
 \draw[step=0.7cm ,gray,very thin] (0,0) grid (18,6);
 % Обводка внешнего контура  \draw[black, thick] (0,0) rectangle (18,6);
 \end{tikzpicture}
\item Найди число, если его $\frac{1}{3}$  часть равна 15. Нарисуй схему.  \\ \nopagebreak \begin{tikzpicture}
 \draw[step=0.7cm ,gray,very thin] (0,0) grid (18,6);
 % Обводка внешнего контура  \draw[black, thick] (0,0) rectangle (18,6);
 \end{tikzpicture}
\item Найди число, если его  14 \%  равна 4200. Нарисуй схему.  \\ \nopagebreak \begin{tikzpicture}
 \draw[step=0.7cm ,gray,very thin] (0,0) grid (18,6);
 % Обводка внешнего контура  \draw[black, thick] (0,0) rectangle (18,6);
 \end{tikzpicture}
\end{enumerate}  
\newpage
\subsection*{Вариант 9 для  Кострицкой Есении Ивановны. Дроби \\ \today   \currenttime} 
  \begin{enumerate}
\item Найди число, если его  20 \%  равна 28000. Нарисуй схему.  \\ \nopagebreak \begin{tikzpicture}
 \draw[step=0.7cm ,gray,very thin] (0,0) grid (18,6);
 % Обводка внешнего контура  \draw[black, thick] (0,0) rectangle (18,6);
 \end{tikzpicture}
\item Найди число, если его $\frac{1}{7}$  часть равна 21. Нарисуй схему.  \\ \nopagebreak \begin{tikzpicture}
 \draw[step=0.7cm ,gray,very thin] (0,0) grid (18,6);
 % Обводка внешнего контура  \draw[black, thick] (0,0) rectangle (18,6);
 \end{tikzpicture}
\item Найди число, если его  8 \%  равна 2400. Нарисуй схему.  \\ \nopagebreak \begin{tikzpicture}
 \draw[step=0.7cm ,gray,very thin] (0,0) grid (18,6);
 % Обводка внешнего контура  \draw[black, thick] (0,0) rectangle (18,6);
 \end{tikzpicture}
\item Найди число, если его $\frac{9}{12}$  часть равна 108. Нарисуй схему.  \\ \nopagebreak \begin{tikzpicture}
 \draw[step=0.7cm ,gray,very thin] (0,0) grid (18,6);
 % Обводка внешнего контура  \draw[black, thick] (0,0) rectangle (18,6);
 \end{tikzpicture}
\item Найди число, если его $\frac{6}{17}$  часть равна 102. Нарисуй схему.  \\ \nopagebreak \begin{tikzpicture}
 \draw[step=0.7cm ,gray,very thin] (0,0) grid (18,6);
 % Обводка внешнего контура  \draw[black, thick] (0,0) rectangle (18,6);
 \end{tikzpicture}
\item Найди число, если его $\frac{1}{6}$  часть равна 30. Нарисуй схему.  \\ \nopagebreak \begin{tikzpicture}
 \draw[step=0.7cm ,gray,very thin] (0,0) grid (18,6);
 % Обводка внешнего контура  \draw[black, thick] (0,0) rectangle (18,6);
 \end{tikzpicture}
\end{enumerate}  
\newpage
\subsection*{Вариант 10 для  Купцова Даниила Андреевича. Дроби \\ \today   \currenttime} 
  \begin{enumerate}
\item Найди число, если его $\frac{2}{7}$  часть равна 14. Нарисуй схему.  \\ \nopagebreak \begin{tikzpicture}
 \draw[step=0.7cm ,gray,very thin] (0,0) grid (18,6);
 % Обводка внешнего контура  \draw[black, thick] (0,0) rectangle (18,6);
 \end{tikzpicture}
\item Найди число, если его $\frac{1}{5}$  часть равна 45. Нарисуй схему.  \\ \nopagebreak \begin{tikzpicture}
 \draw[step=0.7cm ,gray,very thin] (0,0) grid (18,6);
 % Обводка внешнего контура  \draw[black, thick] (0,0) rectangle (18,6);
 \end{tikzpicture}
\item Найди число, если его $\frac{9}{17}$  часть равна 153. Нарисуй схему.  \\ \nopagebreak \begin{tikzpicture}
 \draw[step=0.7cm ,gray,very thin] (0,0) grid (18,6);
 % Обводка внешнего контура  \draw[black, thick] (0,0) rectangle (18,6);
 \end{tikzpicture}
\item Найди число, если его $\frac{1}{3}$  часть равна 24. Нарисуй схему.  \\ \nopagebreak \begin{tikzpicture}
 \draw[step=0.7cm ,gray,very thin] (0,0) grid (18,6);
 % Обводка внешнего контура  \draw[black, thick] (0,0) rectangle (18,6);
 \end{tikzpicture}
\item Найди число, если его  7 \%  равна 11200. Нарисуй схему.  \\ \nopagebreak \begin{tikzpicture}
 \draw[step=0.7cm ,gray,very thin] (0,0) grid (18,6);
 % Обводка внешнего контура  \draw[black, thick] (0,0) rectangle (18,6);
 \end{tikzpicture}
\item Найди число, если его  11 \%  равна 14300. Нарисуй схему.  \\ \nopagebreak \begin{tikzpicture}
 \draw[step=0.7cm ,gray,very thin] (0,0) grid (18,6);
 % Обводка внешнего контура  \draw[black, thick] (0,0) rectangle (18,6);
 \end{tikzpicture}
\end{enumerate}  
\newpage
\subsection*{Вариант 11 для  Лукашовой Мирославы Евгеньевны. Дроби \\ \today   \currenttime} 
  \begin{enumerate}
\item Найди число, если его  10 \%  равна 5000. Нарисуй схему.  \\ \nopagebreak \begin{tikzpicture}
 \draw[step=0.7cm ,gray,very thin] (0,0) grid (18,6);
 % Обводка внешнего контура  \draw[black, thick] (0,0) rectangle (18,6);
 \end{tikzpicture}
\item Найди число, если его  9 \%  равна 18000. Нарисуй схему.  \\ \nopagebreak \begin{tikzpicture}
 \draw[step=0.7cm ,gray,very thin] (0,0) grid (18,6);
 % Обводка внешнего контура  \draw[black, thick] (0,0) rectangle (18,6);
 \end{tikzpicture}
\item Найди число, если его $\frac{4}{8}$  часть равна 32. Нарисуй схему.  \\ \nopagebreak \begin{tikzpicture}
 \draw[step=0.7cm ,gray,very thin] (0,0) grid (18,6);
 % Обводка внешнего контура  \draw[black, thick] (0,0) rectangle (18,6);
 \end{tikzpicture}
\item Найди число, если его $\frac{7}{18}$  часть равна 126. Нарисуй схему.  \\ \nopagebreak \begin{tikzpicture}
 \draw[step=0.7cm ,gray,very thin] (0,0) grid (18,6);
 % Обводка внешнего контура  \draw[black, thick] (0,0) rectangle (18,6);
 \end{tikzpicture}
\item Найди число, если его $\frac{1}{8}$  часть равна 32. Нарисуй схему.  \\ \nopagebreak \begin{tikzpicture}
 \draw[step=0.7cm ,gray,very thin] (0,0) grid (18,6);
 % Обводка внешнего контура  \draw[black, thick] (0,0) rectangle (18,6);
 \end{tikzpicture}
\item Найди число, если его $\frac{1}{6}$  часть равна 60. Нарисуй схему.  \\ \nopagebreak \begin{tikzpicture}
 \draw[step=0.7cm ,gray,very thin] (0,0) grid (18,6);
 % Обводка внешнего контура  \draw[black, thick] (0,0) rectangle (18,6);
 \end{tikzpicture}
\end{enumerate}  
\newpage
\subsection*{Вариант 12 для  Мухаметшиной Алисы Рифатовны. Дроби \\ \today   \currenttime} 
  \begin{enumerate}
\item Найди число, если его $\frac{8}{19}$  часть равна 152. Нарисуй схему.  \\ \nopagebreak \begin{tikzpicture}
 \draw[step=0.7cm ,gray,very thin] (0,0) grid (18,6);
 % Обводка внешнего контура  \draw[black, thick] (0,0) rectangle (18,6);
 \end{tikzpicture}
\item Найди число, если его $\frac{5}{12}$  часть равна 60. Нарисуй схему.  \\ \nopagebreak \begin{tikzpicture}
 \draw[step=0.7cm ,gray,very thin] (0,0) grid (18,6);
 % Обводка внешнего контура  \draw[black, thick] (0,0) rectangle (18,6);
 \end{tikzpicture}
\item Найди число, если его $\frac{1}{8}$  часть равна 72. Нарисуй схему.  \\ \nopagebreak \begin{tikzpicture}
 \draw[step=0.7cm ,gray,very thin] (0,0) grid (18,6);
 % Обводка внешнего контура  \draw[black, thick] (0,0) rectangle (18,6);
 \end{tikzpicture}
\item Найди число, если его  2 \%  равна 2600. Нарисуй схему.  \\ \nopagebreak \begin{tikzpicture}
 \draw[step=0.7cm ,gray,very thin] (0,0) grid (18,6);
 % Обводка внешнего контура  \draw[black, thick] (0,0) rectangle (18,6);
 \end{tikzpicture}
\item Найди число, если его  9 \%  равна 16200. Нарисуй схему.  \\ \nopagebreak \begin{tikzpicture}
 \draw[step=0.7cm ,gray,very thin] (0,0) grid (18,6);
 % Обводка внешнего контура  \draw[black, thick] (0,0) rectangle (18,6);
 \end{tikzpicture}
\item Найди число, если его $\frac{1}{3}$  часть равна 18. Нарисуй схему.  \\ \nopagebreak \begin{tikzpicture}
 \draw[step=0.7cm ,gray,very thin] (0,0) grid (18,6);
 % Обводка внешнего контура  \draw[black, thick] (0,0) rectangle (18,6);
 \end{tikzpicture}
\end{enumerate}  
\newpage
\subsection*{Вариант 13 для  Нижегородова Даниила Кирилловича. Дроби \\ \today   \currenttime} 
  \begin{enumerate}
\item Найди число, если его  16 \%  равна 35200. Нарисуй схему.  \\ \nopagebreak \begin{tikzpicture}
 \draw[step=0.7cm ,gray,very thin] (0,0) grid (18,6);
 % Обводка внешнего контура  \draw[black, thick] (0,0) rectangle (18,6);
 \end{tikzpicture}
\item Найди число, если его $\frac{7}{17}$  часть равна 119. Нарисуй схему.  \\ \nopagebreak \begin{tikzpicture}
 \draw[step=0.7cm ,gray,very thin] (0,0) grid (18,6);
 % Обводка внешнего контура  \draw[black, thick] (0,0) rectangle (18,6);
 \end{tikzpicture}
\item Найди число, если его $\frac{1}{5}$  часть равна 45. Нарисуй схему.  \\ \nopagebreak \begin{tikzpicture}
 \draw[step=0.7cm ,gray,very thin] (0,0) grid (18,6);
 % Обводка внешнего контура  \draw[black, thick] (0,0) rectangle (18,6);
 \end{tikzpicture}
\item Найди число, если его  9 \%  равна 2700. Нарисуй схему.  \\ \nopagebreak \begin{tikzpicture}
 \draw[step=0.7cm ,gray,very thin] (0,0) grid (18,6);
 % Обводка внешнего контура  \draw[black, thick] (0,0) rectangle (18,6);
 \end{tikzpicture}
\item Найди число, если его $\frac{1}{4}$  часть равна 20. Нарисуй схему.  \\ \nopagebreak \begin{tikzpicture}
 \draw[step=0.7cm ,gray,very thin] (0,0) grid (18,6);
 % Обводка внешнего контура  \draw[black, thick] (0,0) rectangle (18,6);
 \end{tikzpicture}
\item Найди число, если его $\frac{6}{15}$  часть равна 90. Нарисуй схему.  \\ \nopagebreak \begin{tikzpicture}
 \draw[step=0.7cm ,gray,very thin] (0,0) grid (18,6);
 % Обводка внешнего контура  \draw[black, thick] (0,0) rectangle (18,6);
 \end{tikzpicture}
\end{enumerate}  
\newpage
\subsection*{Вариант 14 для  Поландова Максима Ивановича. Дроби \\ \today   \currenttime} 
  \begin{enumerate}
\item Найди число, если его  17 \%  равна 28900. Нарисуй схему.  \\ \nopagebreak \begin{tikzpicture}
 \draw[step=0.7cm ,gray,very thin] (0,0) grid (18,6);
 % Обводка внешнего контура  \draw[black, thick] (0,0) rectangle (18,6);
 \end{tikzpicture}
\item Найди число, если его  17 \%  равна 17000. Нарисуй схему.  \\ \nopagebreak \begin{tikzpicture}
 \draw[step=0.7cm ,gray,very thin] (0,0) grid (18,6);
 % Обводка внешнего контура  \draw[black, thick] (0,0) rectangle (18,6);
 \end{tikzpicture}
\item Найди число, если его $\frac{8}{20}$  часть равна 160. Нарисуй схему.  \\ \nopagebreak \begin{tikzpicture}
 \draw[step=0.7cm ,gray,very thin] (0,0) grid (18,6);
 % Обводка внешнего контура  \draw[black, thick] (0,0) rectangle (18,6);
 \end{tikzpicture}
\item Найди число, если его $\frac{2}{9}$  часть равна 18. Нарисуй схему.  \\ \nopagebreak \begin{tikzpicture}
 \draw[step=0.7cm ,gray,very thin] (0,0) grid (18,6);
 % Обводка внешнего контура  \draw[black, thick] (0,0) rectangle (18,6);
 \end{tikzpicture}
\item Найди число, если его $\frac{1}{3}$  часть равна 21. Нарисуй схему.  \\ \nopagebreak \begin{tikzpicture}
 \draw[step=0.7cm ,gray,very thin] (0,0) grid (18,6);
 % Обводка внешнего контура  \draw[black, thick] (0,0) rectangle (18,6);
 \end{tikzpicture}
\item Найди число, если его $\frac{1}{7}$  часть равна 21. Нарисуй схему.  \\ \nopagebreak \begin{tikzpicture}
 \draw[step=0.7cm ,gray,very thin] (0,0) grid (18,6);
 % Обводка внешнего контура  \draw[black, thick] (0,0) rectangle (18,6);
 \end{tikzpicture}
\end{enumerate}  
\newpage
\subsection*{Вариант 15 для  Поповой Ульяны Дмитриевны. Дроби \\ \today   \currenttime} 
  \begin{enumerate}
\item Найди число, если его $\frac{1}{9}$  часть равна 63. Нарисуй схему.  \\ \nopagebreak \begin{tikzpicture}
 \draw[step=0.7cm ,gray,very thin] (0,0) grid (18,6);
 % Обводка внешнего контура  \draw[black, thick] (0,0) rectangle (18,6);
 \end{tikzpicture}
\item Найди число, если его $\frac{8}{11}$  часть равна 88. Нарисуй схему.  \\ \nopagebreak \begin{tikzpicture}
 \draw[step=0.7cm ,gray,very thin] (0,0) grid (18,6);
 % Обводка внешнего контура  \draw[black, thick] (0,0) rectangle (18,6);
 \end{tikzpicture}
\item Найди число, если его  4 \%  равна 6800. Нарисуй схему.  \\ \nopagebreak \begin{tikzpicture}
 \draw[step=0.7cm ,gray,very thin] (0,0) grid (18,6);
 % Обводка внешнего контура  \draw[black, thick] (0,0) rectangle (18,6);
 \end{tikzpicture}
\item Найди число, если его $\frac{7}{18}$  часть равна 126. Нарисуй схему.  \\ \nopagebreak \begin{tikzpicture}
 \draw[step=0.7cm ,gray,very thin] (0,0) grid (18,6);
 % Обводка внешнего контура  \draw[black, thick] (0,0) rectangle (18,6);
 \end{tikzpicture}
\item Найди число, если его $\frac{1}{9}$  часть равна 45. Нарисуй схему.  \\ \nopagebreak \begin{tikzpicture}
 \draw[step=0.7cm ,gray,very thin] (0,0) grid (18,6);
 % Обводка внешнего контура  \draw[black, thick] (0,0) rectangle (18,6);
 \end{tikzpicture}
\item Найди число, если его  10 \%  равна 21000. Нарисуй схему.  \\ \nopagebreak \begin{tikzpicture}
 \draw[step=0.7cm ,gray,very thin] (0,0) grid (18,6);
 % Обводка внешнего контура  \draw[black, thick] (0,0) rectangle (18,6);
 \end{tikzpicture}
\end{enumerate}  
\newpage
\subsection*{Вариант 16 для  Тильдикова Ярослава Александровича. Дроби \\ \today   \currenttime} 
  \begin{enumerate}
\item Найди число, если его  10 \%  равна 5000. Нарисуй схему.  \\ \nopagebreak \begin{tikzpicture}
 \draw[step=0.7cm ,gray,very thin] (0,0) grid (18,6);
 % Обводка внешнего контура  \draw[black, thick] (0,0) rectangle (18,6);
 \end{tikzpicture}
\item Найди число, если его $\frac{1}{6}$  часть равна 42. Нарисуй схему.  \\ \nopagebreak \begin{tikzpicture}
 \draw[step=0.7cm ,gray,very thin] (0,0) grid (18,6);
 % Обводка внешнего контура  \draw[black, thick] (0,0) rectangle (18,6);
 \end{tikzpicture}
\item Найди число, если его  10 \%  равна 7000. Нарисуй схему.  \\ \nopagebreak \begin{tikzpicture}
 \draw[step=0.7cm ,gray,very thin] (0,0) grid (18,6);
 % Обводка внешнего контура  \draw[black, thick] (0,0) rectangle (18,6);
 \end{tikzpicture}
\item Найди число, если его $\frac{5}{11}$  часть равна 55. Нарисуй схему.  \\ \nopagebreak \begin{tikzpicture}
 \draw[step=0.7cm ,gray,very thin] (0,0) grid (18,6);
 % Обводка внешнего контура  \draw[black, thick] (0,0) rectangle (18,6);
 \end{tikzpicture}
\item Найди число, если его $\frac{1}{11}$  часть равна 88. Нарисуй схему.  \\ \nopagebreak \begin{tikzpicture}
 \draw[step=0.7cm ,gray,very thin] (0,0) grid (18,6);
 % Обводка внешнего контура  \draw[black, thick] (0,0) rectangle (18,6);
 \end{tikzpicture}
\item Найди число, если его $\frac{3}{12}$  часть равна 36. Нарисуй схему.  \\ \nopagebreak \begin{tikzpicture}
 \draw[step=0.7cm ,gray,very thin] (0,0) grid (18,6);
 % Обводка внешнего контура  \draw[black, thick] (0,0) rectangle (18,6);
 \end{tikzpicture}
\end{enumerate}  
\newpage
\subsection*{Вариант 17 для  Топчиева Аарона Рустамовича. Дроби \\ \today   \currenttime} 
  \begin{enumerate}
\item Найди число, если его $\frac{9}{14}$  часть равна 126. Нарисуй схему.  \\ \nopagebreak \begin{tikzpicture}
 \draw[step=0.7cm ,gray,very thin] (0,0) grid (18,6);
 % Обводка внешнего контура  \draw[black, thick] (0,0) rectangle (18,6);
 \end{tikzpicture}
\item Найди число, если его $\frac{6}{11}$  часть равна 66. Нарисуй схему.  \\ \nopagebreak \begin{tikzpicture}
 \draw[step=0.7cm ,gray,very thin] (0,0) grid (18,6);
 % Обводка внешнего контура  \draw[black, thick] (0,0) rectangle (18,6);
 \end{tikzpicture}
\item Найди число, если его  10 \%  равна 4000. Нарисуй схему.  \\ \nopagebreak \begin{tikzpicture}
 \draw[step=0.7cm ,gray,very thin] (0,0) grid (18,6);
 % Обводка внешнего контура  \draw[black, thick] (0,0) rectangle (18,6);
 \end{tikzpicture}
\item Найди число, если его $\frac{1}{11}$  часть равна 88. Нарисуй схему.  \\ \nopagebreak \begin{tikzpicture}
 \draw[step=0.7cm ,gray,very thin] (0,0) grid (18,6);
 % Обводка внешнего контура  \draw[black, thick] (0,0) rectangle (18,6);
 \end{tikzpicture}
\item Найди число, если его  9 \%  равна 13500. Нарисуй схему.  \\ \nopagebreak \begin{tikzpicture}
 \draw[step=0.7cm ,gray,very thin] (0,0) grid (18,6);
 % Обводка внешнего контура  \draw[black, thick] (0,0) rectangle (18,6);
 \end{tikzpicture}
\item Найди число, если его $\frac{1}{3}$  часть равна 24. Нарисуй схему.  \\ \nopagebreak \begin{tikzpicture}
 \draw[step=0.7cm ,gray,very thin] (0,0) grid (18,6);
 % Обводка внешнего контура  \draw[black, thick] (0,0) rectangle (18,6);
 \end{tikzpicture}
\end{enumerate}  
\newpage
